% Results section. Numerical claims are restricted to the verified facts in T24.
\section{Results}

\subsection{Descriptive picture}

The national weighted prevalence of severe food insecurity was 29.10\%, while the unweighted prevalence in the analysis sample was 27.74\%. These summaries refer to different estimands. The fitted likelihood uses the unweighted household records, so the unweighted proportion is the directly comparable sample summary; the weighted estimate remains the survey-weighted national result. We report both rather than treating the fitted surface as a replacement for the weighted prevalence. This distinction also applies to the area results below: the model describes probabilities for the unweighted sample quantity, while the weighted figure describes the national survey estimate.

Direct state estimates ranged from 4.6\% in Niger to 85.0\% in Yobe, with a median of 23.5\%. The range shows that the state summaries vary widely, and the median locates the middle of that distribution well below its upper endpoint. These state values provide a national-scale description, but they do not show how estimates vary among local government areas within a state. The LGA-level posterior summaries therefore add a finer geographic description rather than replacing the state estimates.

Affirmation was less common for more severe food experiences. The item-level affirmation rate fell monotonically from 71\% for worrying about food to 33.5\% for not eating for a whole day, with no inversions in the ordering. The monotone gradient is visible across the reported endpoints and was not interrupted by an item with a higher affirmation rate at a more severe point in the sequence. It describes the distribution of reported experiences; it is not an alternative prevalence definition for the area model.

The two spatial-dependence statistics differed by outcome measure. For severe food insecurity, Moran's \(I\) was 0.204, with permutation \(p=0.018\) from 9,999 draws. For the broader moderate-or-severe measure, Moran's \(I\) was \(-0.012\), with \(p=0.427\). Thus, the permutation result indicates spatial dependence for the severe measure, whereas the corresponding test did not detect spatial dependence for the broader measure. The second result is reported alongside the first because the measures do not yield the same spatial pattern.

The non-significant result should not be read as proof that the broader measure has no spatial association. It records that this test did not distinguish its statistic from the null distribution. The significant result for severe food insecurity is also a global summary: it does not imply that all locations share the same level of dependence. The state range and the LGA estimates describe geographic variation, while the two Moran statistics summarize different outcomes at a broader scale.

\subsection{Area estimates}

The posterior means for the 774 LGAs ranged from 0.017 to 0.918. Figure~\ref{fig:estimates} displays those means in sorted order with their intervals, so readers can see both the spread of the fitted estimates and their area-specific uncertainty. The displayed range is a range of model estimates, not a range of direct sample proportions. The intervals are needed to read the ordered means: two nearby point estimates need not have the same posterior uncertainty, and the plot preserves that distinction rather than reducing the distribution to its endpoints.

The endpoints alone do not describe the full distribution of area estimates. They do not show how many LGAs lie near the crisis threshold, nor do they convey how certain the model is about an individual area. The sorted display keeps those two features visible together. It should therefore be read as a collection of posterior area summaries, rather than as a ranking of observed household proportions.

Sample coverage was incomplete at LGA level. A sampled cluster was present in 654 LGAs, while 120 had no sampled cluster; those 120 areas contain 11.04\% of the population. For the unclustered LGAs, the model's posterior means ranged from 0.046 to 0.584, and their mean exceedance probability was 0.132. These are posterior predictions rather than direct observations. No sampled cluster contributes a local household outcome in those areas, so the reported means and exceedance probabilities are model output and should not be described as measured LGA prevalence.

The mean exceedance probability of 0.132 is a summary of posterior probabilities across the unclustered LGAs. It is not the share of households in those areas with severe food insecurity, and it is not an observed prevalence. The posterior mean and the exceedance probability answer different questions: one summarizes the predicted area prevalence, and the other summarizes the model's probability that an area crosses the chosen line. Both quantities for these areas remain predictions because the survey supplies no sampled cluster there.

The model summaries also distinguish the area and cluster variation represented in the fit. The reported LGA precision was 0.904 and the cluster precision was 4.784. The spatial mixing parameter \(\phi\) was 0.971, with a 95\% interval from 0.903 to 0.997. These parameter summaries accompany the area estimates but do not remove the difference between an LGA with observed households and an LGA whose estimate is produced without a sampled cluster.

The parameter summaries and the calibration results describe separate aspects of the fit. The mixing parameter and precisions summarize fitted variation, whereas the holdout results assess prediction against areas excluded from the corresponding fit. A parameter estimate within its reported interval is therefore not, by itself, a validation result for the area-specific predictions.

The model fit took 40--55 seconds, and its WAIC was 37,942.17 across runs. The posterior summaries reported here use 20,000 draws with seed 2026-11-22. These computational summaries identify the fit and draw set behind the reported quantities; they are not evidence that every area-specific interval has the same predictive performance. The calibration results below assess prediction under distinct holdout designs.

\begin{figure}[htbp]
  \centering
  \fbox{\parbox{0.88\linewidth}{\centering Figure output pending.}}
  \caption{Posterior mean prevalence for each of the 774 LGAs, ordered from lowest to highest, with the corresponding interval for each area. Output file: \todo{figure output filename}.}
  \label{fig:estimates}
\end{figure}

\subsection{The count}

At the threshold \(c=0.50\), the posterior expected number of LGAs above the threshold was \(E[C]=114.975\) of 774. The 95\% posterior interval was 101 to 130, with posterior standard deviation 7.383. This interval summarizes the model's joint posterior draws and is conditional on the fitted dependence structure; the marginal exceedance probabilities alone do not determine it. The feasible marginal-only range for the count's standard deviation is reported in the next subsection, where the distinction between the model interval and the information in the marginals is made explicit.

The variance decomposition separates the contribution of each LGA's exceedance probability from the contribution of pairwise joint exceedance. At \(c=0.50\), the independent term was 40.258 and the pairwise term was 14.248, giving a pairwise-term ratio of 0.354. The identity residual was \(-2.7\times 10^{-3}\). The exchangeable-areas comparator had a ratio of 0.187. These summaries show that the pairwise component is positive in the fitted model, while the comparator gives a smaller ratio under its alternative dependence structure.

The pairwise contribution is the part that is absent when the area indicators are treated as independent. Its positive value means that the fitted joint draws contain additional count variation beyond the independent term. The exchangeable-area ratio is a comparator for this same count summary; it should not be read as a decomposition into a separate causal contribution from geography. The reported ratios describe the two fitted dependence structures, not a test of whether one structure is correct.

The pairwise term has visible Monte Carlo variation. The split-half estimates were 13.37 and 15.12, a difference of approximately \(\pm 1.7\) around the reported value. We therefore report the interval and variance decomposition together, rather than treating the pairwise contribution as an exact constant. The split-half comparison concerns numerical variation in this posterior summary; it does not change the separate fact that the joint dependence structure is an assumption of the fitted model.

Posterior classification also leaves a group of areas close to the crisis line. Thirty-two LGAs had exceedance probability above 0.95, 509 were classified below the threshold, and 40 had probabilities between 0.40 and 0.60. The expected count changed with the area-level threshold: \(E[C]\) was 196.28 at \(c=0.40\) and 61.28 at \(c=0.60\), compared with 114.975 at \(c=0.50\). These three expected counts summarize threshold sensitivity; the variance decomposition in Table~\ref{tab:count} is reported at the central threshold.

\begin{table}[htbp]
  \centering
  \caption{Posterior count summaries across thresholds and variance decomposition at \(c=0.50\). The 95\% interval at \(c=0.50\) is conditional on the fitted joint dependence model and is not determined by the marginal exceedance probabilities alone.}
  \label{tab:count}
  \begin{tabular}{lrr}
    \toprule
    Threshold \(c\) & \(E[C]\) & 95\% posterior interval \\
    \midrule
    0.40 & 196.28 & 180--213 \\
    0.50 & 114.975 & 101--130 \\
    0.60 & 61.28 & 51--73 \\
    \midrule
    \multicolumn{3}{l}{Variance decomposition and comparator at \(c=0.50\)} \\
    Posterior standard deviation & \multicolumn{2}{r}{7.383} \\
    Independent variance term & \multicolumn{2}{r}{40.258} \\
    Pairwise variance term & \multicolumn{2}{r}{14.248} \\
    Pairwise-term ratio & \multicolumn{2}{r}{0.354} \\
    Exchangeable-areas comparator ratio & \multicolumn{2}{r}{0.187} \\
    Identity residual & \multicolumn{2}{r}{\(-2.7\times 10^{-3}\)} \\
    \bottomrule
  \end{tabular}
\end{table}

\subsection{What the marginals can and cannot pin down}

The marginal exceedance probabilities do not identify the variance of the count. By Proposition~\ref{prop:sharp}, for the 774 supplied marginals, the feasible range of \(\operatorname{Var}(C)\) is 0.0244 to 5492.63, corresponding to a standard-deviation range of 0.156 to 74.11. The lower and upper endpoints are attained by the cyclic and comonotone couplings, respectively. The fitted model's variance of 54.509 lies 0.99\% of the way up this feasible range. Thus, the \(E[C]\) value is fixed by the marginals, but the uncertainty around that count depends on how the exceedances are joined across areas.

The reported 95\% count interval of 101 to 130 is therefore a property of the fitted model's joint dependence structure, not a property of the data's marginal probabilities alone. The same marginals permit a count standard deviation anywhere from 0.156 to 74.11, so they cannot by themselves select the reported interval. The distinction is also visible in the tail probability: under the comonotone coupling, \(\Pr(C\geq 101)=0.476\); under the cyclic coupling it is 1.000; under the fitted model it is 0.984. These values use the same marginals but different joint couplings. The model interval must consequently be read as conditional on the dependence structure estimated by the model, not as an interval that the marginal data alone establish.

\begin{figure}[htbp]
  \centering
  \fbox{\parbox{0.88\linewidth}{\centering Figure output pending.}}
  \caption{The feasible range of \(\operatorname{Var}(C)\) given the area-specific marginal exceedance probabilities, with the fitted model's variance marked. Output file: \todo{figure output filename}.}
  \label{fig:feasible}
\end{figure}

\subsection{Partial identification}

Allowing the unclustered areas to contribute either zero or one to the threshold count gives an identified interval from 87 to 231 LGAs, with width 144. The observed-only count was 99.13, with an interval from 87 to 111. The full identified interval contains the model-based 101 to 130 interval, but the latter is conditional on the fitted joint dependence structure and is not determined by the marginals alone. Of the identified interval's width, 120 counts arise from the unrestricted contribution of the unclustered areas. The lower endpoint still uses model-smoothed prevalences for LGAs with sampled clusters, so the interval combines uncertainty about the unsampled areas with the fitted estimates for the observed areas.

The population-share bounds tell a different story from the count of areas because the unclustered LGAs do not contain an equal share of the population. The observed-only share of people living in a crisis LGA was 13.72\%, and the identified interval for the national share was 10.63\% to 24.90\%. The width of the population-share bound was 0.11037, equal to \(1-p\). This width describes the share of people whose crisis-area status is not fixed by the observed areas; it is not the fraction of geographic units without a cluster.

The count interval and the population-share interval should not be substituted for one another. The count gives every LGA one unit of weight, while the population share gives each area weight according to the people represented there. The unclustered areas account for a smaller share of people than their share of geographic units, so the population-share bounds are narrower than an area count alone would suggest. The reported endpoints keep the uncertainty attached to people distinct from the uncertainty attached to the number of places.

Under the monotonicity variant, the population-share interval was 0.1372 to 0.2324, a reduction in width of 1.3 percentage points. This variant narrows the interval relative to the unrestricted bound, but it relies on an additional ordering assumption. We report it as a sensitivity result rather than replacing the unrestricted interval. The count and population-share results are posterior intervals for sample-analogue identified sets, not frequentist confidence sets.

\begin{figure}[htbp]
  \centering
  \fbox{\parbox{0.88\linewidth}{\centering Figure output pending.}}
  \caption{The identified interval for the count compared with the model-based count interval. The model interval is conditional on the fitted joint dependence structure and is not identified by the marginal probabilities alone. Output file: \todo{figure output filename}.}
  \label{fig:identified}
\end{figure}

\subsection{Calibration}

The four calibration designs gave different results depending on whether prediction concerned a single area among observed neighbours or a whole region without its local observations. Table~\ref{tab:calibration} reports exposure, the mean and standard deviation of \(z\), coverage, and the result against each pre-declared criterion. The pass rule required mean \(z\) between \(-0.15\) and 0.15, standard deviation between 0.85 and 1.15, and coverage between 0.90 and 0.98. Each area-level design held out 200 LGAs, except the whole-region design; the area-level designs used groups of 40 for prediction.

\begin{table}[htbp]
  \centering
  \caption{Spatial calibration results. Exposure is the mean number of a held-out LGA's neighbours that were not seen by the fit used to predict it. The pre-declared pass bands were mean \(z\) in \([-0.15,0.15]\), standard deviation of \(z\) in \([0.85,1.15]\), and coverage in \([0.90,0.98]\).}
  \label{tab:calibration}
  \begin{tabular}{lrrrrl}
    \toprule
    Design & Exposure & Mean \(z\) & SD \(z\) & Coverage & Verdicts (mean/SD/coverage) \\
    \midrule
    Whole region removed & 5.4 & \(-0.654\) & 1.648 & 0.775 & Fail / Fail / Fail \\
    Single areas, neighbours kept & 0.88 & \(-0.052\) & 1.111 & 0.940 & Pass / Pass / Pass \\
    Matched over-exposed & 1.67 & \(-0.068\) & 1.181 & 0.900 & Pass / Fail / Pass \\
    Exact exposure of the 120 & 0.95 & \(-0.058\) & 1.178 & 0.890 & Pass / Fail / Fail \\
    \bottomrule
  \end{tabular}
\end{table}

The whole-region design failed all three pre-declared criteria. Its mean \(z\) of \(-0.654\) fell outside the allowed range, its standard deviation of 1.648 exceeded the upper limit, and its coverage of 0.775 fell below the lower limit. The failure is therefore not confined to coverage: the mean error and the spread of standardized errors also failed their respective checks. This design tests prediction of an area in a region that was wholly withheld, and the results do not support treating the resulting regional predictions as calibrated.

The single-area design met all three criteria, with mean \(z=-0.052\), standard deviation 1.111, and coverage 0.940. The matched over-exposed design passed the mean and coverage checks but failed the standard-deviation check at 1.181. The exact-exposure design, which matches the exposure of the 120 unclustered LGAs, passed the mean check but failed both the standard-deviation and coverage checks: its standard deviation was 1.178 and its coverage was 0.890. The exact-exposure result is therefore not a full calibration pass, even though its mean standardized error was within the pre-declared range.

The three area-level designs with exposure between 0.88 and 1.67 had coverages of 0.940, 0.890, and 0.900. These differences are within the noise of the held-out groups and do not support a monotone relationship between exposure and coverage. The supported contrast is between predicting a single area while its neighbours remain in the fit and predicting a whole region after its neighbours are removed. The whole-region coverage was 0.775, below the results for the three area-level designs, but the middle values should not be presented as a coverage curve.

At the exposure corresponding to the unclustered areas, group-level coverage ranged from 0.878 to 0.900. An interval inflation factor of \(k=1.178\) raised coverage to 0.940 after widening. The out-of-fold coverage values for the factor were 0.875, 1.000, 0.925, 0.975, and 0.950. These results report the behavior of intervals for areas under the stated prediction designs; they do not establish calibration for the joint count distribution.

The pooled coverage after widening is within the pre-declared coverage band, but the out-of-fold values vary across prediction groups. The exact-exposure design before widening still failed its coverage criterion, so the unadjusted result remains a failure. The widening result is a separate summary for area intervals and does not turn the original exact-exposure check into a pass. Neither result supplies a calibration test for the count formed from jointly dependent area indicators.

The correlation with direct estimates was 0.974 in sample, 0.768 out of sample under the single-area design, and 0.195 when a whole region was withheld. The in-sample and single-area results describe agreement when neighbouring areas remain represented in the fit, while the whole-region result describes a more demanding regional prediction. The gap between the latter two is consistent with the contrast in coverage, but neither area-level correlation validates the model-based count interval. That count interval remains conditional on the model's joint dependence structure, which the marginal probabilities alone do not identify.

\begin{figure}[htbp]
  \centering
  \fbox{\parbox{0.88\linewidth}{\centering Figure output pending.}}
  \caption{Observed coverage plotted against nominal coverage for the four spatial holdout designs. Output file: \todo{figure output filename}.}
  \label{fig:calibration}
\end{figure}

\subsection{Area covariates}

Three of the eight candidate area covariates could not be interpolated adequately. Travel time had an interpolation correlation of 0.614, livestock 0.535, and nightlights 0.522; all three were dropped as un-interpolatable. Five were retained. Population density was not interpolated at all: it was computed from population projections and polygon area. The other four were aridity at 0.999, elevation at 0.977, malaria at 0.923 and vegetation at 0.920. Precipitation interpolated just as well, at 0.999, but was dropped because it correlates with aridity at 0.985 and the selection rule keeps one of any such pair. Table~\ref{tab:covariates} reports all candidate correlations and their disposition. The distinction is between variables that could be carried to the area level and variables that could not, not between covariates with and without an association with the outcome.

\begin{table}[htbp]
  \centering
  \caption{Cross-validated interpolation correlations and disposition of the eight candidate area covariates.}
  \label{tab:covariates}
  \begin{tabular}{lrl}
    \toprule
    Candidate & Interpolation correlation & Disposition \\
    \midrule
    Population density & not interpolated & Kept \\
    Aridity & 0.999 & Kept \\
    Precipitation & 0.999 & Dropped: collinear with aridity, $r=0.985$ \\
    Elevation & 0.977 & Kept \\
    Malaria & 0.923 & Kept \\
    Vegetation & 0.920 & Kept \\
    Travel time & 0.614 & Dropped: un-interpolatable \\
    Livestock & 0.535 & Dropped: un-interpolatable \\
    Nightlights & 0.522 & Dropped: un-interpolatable \\
    \bottomrule
  \end{tabular}
\end{table}

The five interpolatable covariates were largely inert in the area model. After they were added, 19 of 20 coefficient intervals contained zero. Regional-design coverage increased from 0.775 to 0.797, but remained below the pre-declared lower limit of 0.90. The covariate refit therefore increased regional coverage without resolving the regional calibration failure, and it did not yield clear nonzero coefficient estimates for most included terms.
